\documentclass[a4paper,10pt]{article}
\usepackage[francais]{babel}
\usepackage{url}

\selectlanguage{francais}
\input{letter-icube-uds.sty}

\selectlanguage{francais}
\begin{document}


\name{St\'ephane \textsc{Genaud}}
\fonction{Ma\^itre de conf\'erences (HDR)}
\telephoneillkirch{+33 (0)3688 54542}
\email{genaud@unistra.fr}

\begin{letter}[Recommendation de Rajni Aron]{to}{INRIA Nancy- Grand Est}


Madame  Rajni ARON  a obtenu  son doctorat  en informatique  de l'Université  de
Thapar,  région  du Penjab,  Inde  en  avril 2013.  Elle  a  obtenu un  MSc.  in
Mathematics \&  Computing de cette même  université en juin 2009.  Elle a obtenu
son diplôme de licence (graduation) en 2007 de l'Université du Penjab.

Au cours de sa thèse, Madame Aron a  eu un très bon niveau de publication de ses
travaux: 3 articles dans des journaux  internationaux de très bon niveau (Future
Generation of Computer Systems, Elsevier + Journal of Supercomputing, Springer +
Journal  of Grid  Computing, Springer)  plus un  journal moins  reconnu. Elle  a
également plusieurs  conférences orales dont  une publiée  dans LNCS et  une par
IEEE.

Le travail de thèse de Madame  Aron concorde parfaitement avec le sujet proposé.
Elle  possède un  expérience  sur la  conception de  middleware  pour gérer  des
ressources de grilles -- expérience également acquise précédemment dans l'équipe
AlGorille -- mais a déjà abordé  de manière pratique les environnements de cloud
tels  ceux  visés dans  le  sujet  de postdoc  (avec  une  certification sur  la
plateforme Aneka).  Elle a donc un profil idéal pour aborder le sujet de postdoc
qui concerne la gestion des ressources  en environnements de type cloud. Le fait
qu'elle puisse  être opérationnelle très rapidement  sur ce postdoc d'un  an est
très important à notre sens.


Notre décision de soutenir la candidature de Madame Aron fait suite à un échange
de mails associé à une demande de  travail dans un temps limité.  Nous lui avons
demandé  de faire  par  écrit, du  vendredi  pour le  lundi  matin, une  analyse
critique  d'un   de  nos  articles  et   de  proposer  des  pistes   futures  de
recherche.  Nous  avons ensuite  procédé  à  une interview  skype.  L'impression
générale est  que Madame Aron  fait preuve d'un  grand dynamisme et  nous sommes
persuadés qu'elle saura faire preuve de beaucoup d'autonomie.

Elle n'a pas de  contrainte familiale ou de calendrier. Elle  a en parallèle une
autre candidature  sur un postdoc à  l'université de Tartu en  Estonie, mais dit
privilégier l'INRIA. 





\end{letter}
\begin{flushright}
%\includegraphics[width=.30\textwidth]{signgenaud.jpg}
\end{flushright}
%\newpage
%\bibliographystyle{plain}
%\bibliography{julien_blaise_ED}


\end{document}
